% Theory section of the revised manuscript (stage S6). Compiles on its own for checking;
% the body between the markers is \input into the manuscript. The statements are checked
% numerically by tests/theory/ (table at the end). Revised on 10 November 2026
% after an independent referee report (docs/FINDINGS_V1.md, F18).
\documentclass[11pt]{article}
\usepackage[a4paper,margin=2.4cm]{geometry}
\usepackage{amsmath,amssymb,amsthm,booktabs}
\newtheorem{theorem}{Theorem}
\newtheorem{lemma}{Lemma}
\newtheorem{proposition}{Proposition}
\newtheorem{corollary}{Corollary}
\theoremstyle{definition}
\newtheorem{definition}{Definition}
\newtheorem{assumption}{Assumption}
\newtheorem{remark}{Remark}
\newcommand{\E}{\mathbb{E}}
\newcommand{\Var}{\operatorname{Var}}
\newcommand{\Cov}{\operatorname{Cov}}
\newcommand{\muH}{\mu_H}
\begin{document}
%%% BEGIN BODY

\section{Problem setup, identification and inference}\label{sec:theory}

\subsection{Setup}
A \emph{response set} is a finite collection $U=\{1,\dots,N\}$ of records. Record $i$
carries an outcome $y_i$, covariates $x_i$ that include a detector score with flag
$d_i\in\{0,1\}$, and an origin $a_i$, the share of the response that is machine-written
($a_i\in\{0,1\}$ unless stated otherwise; $a_i\in[0,1]$ for hybrid authorship). Write
$h_i=1-a_i$ for the human share, $N_H=\sum_U h_i>0$, $N_A=N-N_H$ and $\varepsilon=N_A/N$.
The estimand is the human mean of the response set,
\begin{equation}\label{eq:estimand}
\muH=\frac{\sum_{i\in U}h_i\,y_i}{\sum_{i\in U}h_i},
\end{equation}
and $\mu_A$ is the corresponding machine mean when $N_A>0$. Outcomes, covariates and flags
are observed for every record; origins are not. An \emph{audit} verifies the origin of a
subset $S\subset U$ of $m$ records. Strata are indexed by $g$, with $N_g$ records of which
$m_g$ are audited.

\begin{assumption}[D1, probability audit]\label{a:d1}
$S$ is drawn by a randomised design chosen by the analyst, with inclusion probabilities
$\rho_i=\Pr(i\in S)>0$ for every $i\in U$, known for the audited records. (Designs with
several stages are treated in Lemma~\ref{l:twostage} with stage-wise probabilities.)
\end{assumption}
\begin{assumption}[D2, correct verification]\label{a:d2}
For $i\in S$ the audit returns $a_i$. (Proposition~\ref{p:flip} relaxes this.)
\end{assumption}
\begin{assumption}[D3, complete first phase]\label{a:d3}
$(y_i,x_i,d_i)$ is observed for every $i\in U$.
\end{assumption}
\begin{assumption}[M1, non-differential flag]\label{a:m1}
Within each origin class the flag is uncorrelated with the outcome:
$\Cov(D,Y\mid A=a)=0$ for $a\in\{0,1\}$.
\end{assumption}
D1--D3 are properties of the audit, which the analyst controls. M1 is a property of the
detector and the outcome; without verified origins it can sometimes be refuted but never
verified. It was the identifying condition of the first version of this paper. Below it
is needed only for the moment inversion (Proposition~\ref{p:moment}); the estimators we
recommend do not use it.

In Section~\ref{sec:bias} capital letters $(Y,A,D)$ denote a draw from any joint
distribution of outcome, binary origin and flag, in particular the empirical distribution
of a response set, in which case $\Pr(A{=}1)=\varepsilon$, $\E[Y\mid A{=}0]=\muH$ and the
identities hold exactly on that response set. Write
$\bar f=\Pr(D{=}1\mid A{=}0)$ for the false-positive rate, $\bar s=\Pr(D{=}1\mid A{=}1)$ for
the sensitivity and $C_a=\Cov(D,Y\mid A=a)$. A product of zero and an undefined conditional
moment is read as zero.

\subsection{Why deletion and moment inversion are biased}\label{sec:bias}
The naive mean has bias $\varepsilon(\mu_A-\muH)$. Deleting flagged records gives
$\E[Y\mid D{=}0]$.

\begin{lemma}[Deletion bias]\label{l:del}
Let $\bar f<1$ and $\varepsilon_{\mathrm{res}}=\Pr(A{=}1\mid D{=}0)
=\varepsilon(1-\bar s)/\{(1-\varepsilon)(1-\bar f)+\varepsilon(1-\bar s)\}$. Then
\begin{equation}\label{eq:del}
\E[Y\mid D{=}0]-\muH
=\varepsilon_{\mathrm{res}}\,\Delta_{\mathrm{surv}}-\frac{C_0}{1-\bar f},
\qquad
\Delta_{\mathrm{surv}}=\E[Y\mid A{=}1,D{=}0]-\E[Y\mid A{=}0,D{=}0].
\end{equation}
\end{lemma}
\begin{proof}
By total probability over the origin of the unflagged records,
$\E[Y\mid D{=}0]=\E[Y\mid A{=}0,D{=}0]+\varepsilon_{\mathrm{res}}\Delta_{\mathrm{surv}}$.
Among humans, $\E[Y(1-D)\mid A{=}0]=\muH(1-\bar f)-C_0$, so
$\E[Y\mid A{=}0,D{=}0]=\muH-C_0/(1-\bar f)$.
\end{proof}
The first term is residual contamination: machine records that survive deletion, times
their outcome gap to the surviving humans. The second is a selection effect on the humans
themselves: if falsely flagged people answer differently from other people ($C_0\neq0$),
deleting them moves the mean even when the detector catches every machine record. With a
style stratum $G$ and $f(G)=\Pr(D{=}1\mid A{=}0,G)$, $C_0=\Cov(f(G),\E[Y\mid A{=}0,G])
+\E[\Cov(D,Y\mid A{=}0,G)]$; the first submission assumed the last term away.

\begin{proposition}[Moment inversion under M1]\label{p:moment}
Let $0<\varepsilon<1$ and $\bar s\neq\bar f$, with $(\bar s,\bar f,\varepsilon)$ known.
The system
\begin{align}\label{eq:system}
\E[YD]&=(1-\varepsilon)\bar f\,\theta_H+\varepsilon\bar s\,\theta_A,&
\E[Y(1-D)]&=(1-\varepsilon)(1-\bar f)\,\theta_H+\varepsilon(1-\bar s)\,\theta_A
\end{align}
in the unknowns $(\theta_H,\theta_A)$ has a unique solution $(\tilde\mu_H,\tilde\mu_A)$, and under M1
$\tilde\mu_H=\muH$.
\end{proposition}
\begin{proposition}[Bias of the moment inversion without M1]\label{p:momentbias}
Under the conditions of Proposition~\ref{p:moment} but without M1,
\begin{equation}\label{eq:moment}
\tilde\mu_H=\muH-\frac{C_0+\frac{\varepsilon}{1-\varepsilon}\,C_1}{\bar s-\bar f}.
\end{equation}
\end{proposition}
\begin{proof}[Proof of Propositions~\ref{p:moment} and \ref{p:momentbias}]
The determinant of \eqref{eq:system} is $\varepsilon(1-\varepsilon)(\bar f-\bar s)\neq0$ and
Cramer's rule gives $\tilde\mu_H=\{(1-\bar s)\E[YD]-\bar s\,\E[Y(1-D)]\}/
\{(1-\varepsilon)(\bar f-\bar s)\}$. In truth
$\E[YD]=(1-\varepsilon)(\bar f\muH+C_0)+\varepsilon(\bar s\mu_A+C_1)$ and
$\E[Y(1-D)]=(1-\varepsilon)\{(1-\bar f)\muH-C_0\}+\varepsilon\{(1-\bar s)\mu_A-C_1\}$.
Substituting, the terms in $\mu_A$ cancel and the numerator equals
$(1-\varepsilon)\{(\bar f-\bar s)\muH+C_0\}+\varepsilon C_1$, which gives
\eqref{eq:moment}; under M1, $C_0=C_1=0$.
\end{proof}
The inversion is therefore biased whenever $C_0+\frac{\varepsilon}{1-\varepsilon}C_1\neq0$,
and for a given violation the bias is inversely proportional to the detector's separation
$\bar s-\bar f$, the analogue of a weak-instrument bias.

\subsection{Design-based identification}\label{sec:ident}
For a function $z$ on $U$ write $t_z=\sum_U z_i$ and $\hat t_z=\sum_S z_i/\rho_i$ for its
Horvitz--Thompson estimator; under D1, $\E_p[\hat t_z]=t_z$ for every $z$ that does not
depend on $S$, where $\E_p$ is expectation over the audit design with the response set held
fixed.

\begin{definition}[Difference-ratio family]\label{d:family}
Let $\eta:U\to\mathbb{R}$ be a predictor of the origin, computed from $(x_i,d_i)$, and
$\lambda\in\mathbb{R}$. Put
\begin{equation}\label{eq:family}
\hat T_\lambda=\lambda\sum_{i\in U}y_i(1-\eta_i)+\sum_{i\in S}\frac{y_i\{h_i-\lambda(1-\eta_i)\}}{\rho_i},
\qquad
\hat H_\lambda=\lambda\sum_{i\in U}(1-\eta_i)+\sum_{i\in S}\frac{h_i-\lambda(1-\eta_i)}{\rho_i},
\end{equation}
and $\hat\mu=\hat T_\lambda/\hat H_\lambda$, reported when $\hat H_\lambda>0$ (for a tuned member, also at $\lambda=1$) and the audit
contains at least two human records.
\end{definition}
The family contains the two-phase H\'ajek estimator ($\lambda=0$), flag-PPI
($\eta_i=d_i$, $\lambda=1$), RC-PPI ($\eta_i$ a fitted probability that record $i$ is
machine-written, $\lambda=1$) and their power-tuned versions (flag-PPI++, RC-PPI++; a
data-chosen $\lambda$).

\begin{theorem}[Identification without M1]\label{t:ident}
Assume D1--D3, with origins in $[0,1]$, and let $\eta$ and $\lambda$ not depend on $S$. Then
\begin{enumerate}
\item $\E_p[\hat T_\lambda]=N_H\,\muH$ and $\E_p[\hat H_\lambda]=N_H$;
\item whenever $\hat H_\lambda\neq0$,
\begin{equation}\label{eq:repr}
\hat\mu-\muH=\frac{\hat t_e-t_e}{\hat H_\lambda},
\qquad e_i=\{h_i-\lambda(1-\eta_i)\}(y_i-\muH).
\end{equation}
\end{enumerate}
No assumption links the predictor, the flag or the outcome: M1 is not used and $\eta$ may
be arbitrarily wrong.
\end{theorem}
\begin{proof}
(1) The sums over $S$ are Horvitz--Thompson estimators of $\sum_U y_i\{h_i-\lambda(1-\eta_i)\}$
and $\sum_U\{h_i-\lambda(1-\eta_i)\}$; adding the first-phase sums leaves $\sum_U y_ih_i=N_H\muH$
and $\sum_U h_i=N_H$. (2) $\hat T_\lambda-\muH\hat H_\lambda=\lambda\sum_U(y_i-\muH)(1-\eta_i)+\hat t_e$,
and $t_e=\sum_U h_i(y_i-\muH)-\lambda\sum_U(1-\eta_i)(y_i-\muH)=-\lambda\sum_U(1-\eta_i)(y_i-\muH)$
because $\sum_U h_i(y_i-\muH)=0$ by \eqref{eq:estimand}.
\end{proof}
The human mean is thus identified by the audit design alone: numerator and denominator
are design-unbiased whatever the detector does, and the error of the ratio is the
Horvitz--Thompson error of the residual $e$. The detector enters only through $e$. For
$\lambda=1$ the residual is $(\eta_i-a_i)(y_i-\muH)$, the prediction error times the outcome
deviation, which vanishes for a perfect predictor. Two qualifications apply. First,
$\hat\mu$ is a ratio and has the usual ratio bias of order $1/m$, which is not zero in
small audits. Second, the theorem is unconditional: an estimator that is reported only
when the audit contains enough records of each class (as RC-PPI must be, since its
predictor is fitted on the audit) is conditioned on an event that depends on the origins,
and the theorem does not cover that conditional distribution. For this reason the
primary (intention-to-treat) analyses score a complete rule (RC-PPI++ when feasible,
flag-PPI on the same audit otherwise) over all replications; analyses restricted to
feasible replications are conditional and are not covered by Theorem~\ref{t:ident}.

\subsection{Variance, intervals and efficiency}\label{sec:variance}
Let $u_i=h_i(y_i-\muH)$ and $v_i(\eta)=(1-\eta_i)(y_i-\muH)$, so that
$e=u-\lambda v(\eta)$ and $t_u=0$.

\begin{theorem}[Linearised variance and normal limit]\label{t:clt}
Assume D1--D3 and consider a sequence of response sets with single-stage stratified simple
random sampling without replacement, $m\to\infty$ and $m_g\ge2$, along which
(i) $N_H/N$ is bounded away from zero, $y$ and all predictions are bounded, and no stratum
dominates the variance: $m\,N^{-2}\sum_gN_g^2/m_g$ is bounded;
(ii) there are fixed $\eta^\ast$ and $\lambda^\ast$ such that, with
$e^\ast=u-\lambda^\ast v(\eta^\ast)$, the Horvitz--Thompson total $\hat t_{e^\ast}$ satisfies
a central limit theorem under the design, the stratified variance estimator
$\hat V(e^\ast)$ is ratio-consistent for its variance, and the limit is not degenerate:
$\liminf m\,N^{-2}\Var_p(\hat t_{e^\ast})>0$;
(iii) the fitted predictor and tuning constant are stable: each fitted model $\hat\eta$
satisfies $N^{-1}\sum_U(\hat\eta_i-\eta^\ast_i)^2\to0$, the predictions $\tilde\eta$ used
for the audited records satisfy $N^{-1}\sum_S(\tilde\eta_i-\eta^\ast_i)^2/\rho_i\to0$, and
$\hat\lambda\to\lambda^\ast$, all in probability;
(iv) every prediction used for an audited record, in both sums of \eqref{eq:family}, comes
from a model fitted without that record (cross-fitting), and
$N^{-1}m^{1/2}\{\hat t_{v(\tilde\eta)-v(\eta^\ast)}-t_{v(\bar\eta)-v(\eta^\ast)}\}\to0$ in
probability, where $\bar\eta$ denotes the predictions used in the first-phase sum
($\bar\eta_i=\tilde\eta_i$ for $i\in S$).
Then $(\hat\mu-\muH)/\widehat{\mathrm{se}}\to\mathcal N(0,1)$ in distribution, with
\begin{equation}\label{eq:se}
\widehat{\mathrm{se}}^{\,2}=\frac{\hat V(\hat e)}{\hat H_\lambda^{2}},
\qquad \hat e_i=\{h_i-\hat\lambda(1-\tilde\eta_i)\}(y_i-\hat\mu),
\qquad \hat V(\hat e)=\sum_g N_g^2\Big(1-\frac{m_g}{N_g}\Big)\frac{s^2_{\hat e,g}}{m_g}.
\end{equation}
\end{theorem}
\begin{proof}[Proof sketch]
By (i), the Horvitz--Thompson total of any bounded variable has design variance at most a
constant times $\sum_gN_g^2/m_g=O(N^2/m)$; hence $\hat H_\lambda/N_H\to1$ in probability and
$\hat t_{v}-t_{v}=O_p(Nm^{-1/2})$ for bounded $v$. For fixed
$(\eta,\lambda)=(\eta^\ast,\lambda^\ast)$, \eqref{eq:repr} is exact and the claim follows
from the central limit theorem for $\hat t_{e^\ast}$ and Slutsky's lemma; replacing $\muH$
by $\hat\mu$ in the residual changes $\hat V$ by a factor $1+o_p(1)$, again by (i) and the
non-degeneracy in (ii). For fitted quantities,
$\hat T_\lambda-\muH\hat H_\lambda=\hat t_u-\hat\lambda\{\hat t_{v(\tilde\eta)}-t_{v(\bar\eta)}\}$.
The first term does not involve the predictor. By (iv) the bracket equals
$\hat t_{v(\eta^\ast)}-t_{v(\eta^\ast)}+o_p(Nm^{-1/2})$, and since that difference is
$O_p(Nm^{-1/2})$, replacing $\hat\lambda$ by $\lambda^\ast$ costs $o_p(Nm^{-1/2})$ by (iii).
Consistency of the variance estimator at the fitted residual follows from (i) and (iii).
\end{proof}
Conditions (ii) and (iv) are assumptions, not consequences of (iii). Condition (ii) holds
if $e^\ast$ satisfies a Lindeberg-type condition in the strata that are not fully
enumerated (Bickel and Freedman, 1984, for stratified sampling without replacement); it
fails in the degenerate case of a perfect predictor ($e^\ast\equiv0$), and it is strained
when the residual is non-zero for only a handful of records, the situation
Remark~\ref{r:interval} addresses. The allocation condition in (i) holds under
proportional allocation, where $m\,N^{-2}\sum_gN_g^2/m_g=1$; without it, a stratum with a
fixed number of audited records and a non-vanishing share of the response set keeps
$\hat H_\lambda/N_H$ random in the limit, and the limit distribution is not normal.
Condition (iv) is the cross-fitting requirement of prediction-powered inference
with a fitted predictor (Zrnic and Cand\`es, 2024; Angelopoulos et al., 2023, for the
tuning constant), written for the present design; we assume it rather than derive it, and
regard it as plausible for a low-dimensional parametric model such as the two-covariate
logistic model used here. The part about ``both sums'' matters: a predictor that
interpolates its training labels, applied in the first-phase sum to the audited records it
was trained on, violates (iv) and biases the estimate although (iii) holds. Our
implementation therefore uses, for an audited record, its out-of-fold prediction in both
sums, and the average of the fold models for unaudited records. Folds are balanced within
design strata and origin classes. We do not claim the theorem for the two-stage or
systematic designs that are compared as secondary analyses. The theorem is asymptotic;
with 18 to 90 audited records the coverage of the reported intervals is an empirical
question, answered by endpoint PE2.

\begin{corollary}[Power tuning is asymptotically no worse than the audit-only estimator]\label{c:tune}
For fixed $\eta$ with $V_v=\Var_p(\hat t_{v})>0$, the design variance of
$\hat t_u-\lambda\hat t_v$ is $V_u-2\lambda C_{uv}+\lambda^2V_v$, minimised at
$\lambda^\ast=C_{uv}/V_v$ with value $V_u(1-\varrho_{uv}^2)$; for $\lambda^\ast$ clipped to
$[0,1]$ the value is still at most $V_u$, the variance at $\lambda=0$. Hence, under the
conditions of Theorem~\ref{t:clt}, a power-tuned member of the family has asymptotic
variance no larger than the two-phase H\'ajek estimator, for any predictor.
\end{corollary}
\begin{proof}
A convex quadratic in $\lambda$ restricted to $[0,1]$ attains its minimum at the clipped
vertex, and $0\in[0,1]$.
\end{proof}
The statement is asymptotic: in a small audit the tuning constant is itself estimated,
and with an uninformative predictor the tuned estimator can be somewhat worse than the
H\'ajek estimator. No ordering holds between RC-PPI++ and flag-PPI, which use different
predictors; their comparison is empirical (endpoint PE1). In the implementation
$\hat\lambda=\hat C_{uv}/\hat V_v$ is computed from the audit with $\muH$ replaced by the
$\lambda=1$ estimate, clipped to $[0,1]$, and set to 1 when $\hat V_v=0$.

\begin{remark}[Finite-sample interval]\label{r:interval}
The residuals of a good predictor are sparse: $\hat e_i$ is close to zero unless the
predictor errs on record $i$. In a stratum whose few audited records are all predicted
correctly, the sample variance $s^2_{\hat e,g}$ is near zero although the stratum may
contain errors. We therefore (a) shrink each stratum variance towards the variance
anticipated under a working model in which the binary origins of the audited records of
the stratum are independent Bernoulli($p_g$) variables. Writing
$\hat e_i=z^0_i-a_i\delta_i$ with $\delta_i=y_i-\hat\mu$ and treating $z^0$ and $\delta$ as
fixed,
\begin{equation}\label{eq:anticipated}
\E_\xi\big[s^2_{\hat e,g}\big]=p_g(1-p_g)\,\overline{\delta^2}_g+s^2_{z^0-p_g\delta,\,g},
\end{equation}
which follows from $\E\sum_i(\hat e_i-\bar{\hat e})^2=(1-\tfrac1{m_g})\sum_i\delta_i^2p_g(1-p_g)
+\sum_i(c_i-\bar c)^2$ with $c_i=z^0_i-p_g\delta_i$. We evaluate \eqref{eq:anticipated} at the
Jeffreys estimate $p_g=(k_g+\tfrac12)/(m_g+1)$, $k_g$ being the number of machine records
among the $m_g$ audited, and use
$\{(m_g-1)s^2_{\hat e,g}+\kappa\,\E_\xi[s^2_{\hat e,g}]\}/(m_g-1+\kappa)$ with $\kappa=1$.
(b) We use a Student quantile with Satterthwaite degrees of freedom
$(\sum_g\nu_g)^2\big/\sum_g\{\nu_g^2/(m_g-1)\}$, at least 1, where $\nu_g$ is the
contribution of stratum $g$ to $\hat V$. Strata with a single audited record are merged
with a neighbouring stratum before the variance is computed. This rule is a heuristic
selected on development data; its justification is its measured coverage.
\end{remark}

\subsection{Audit design}\label{sec:design-theory}
RC-PPI needs audited records of both classes to fit its predictor. Let $K$ be the number
of audited machine records.

\begin{proposition}[Audit starvation]\label{p:starve}
Let origins be binary. Under simple random sampling of $m$ records, $K$ is hypergeometric
with parameters $(N,N_A,m)$ and mean $m\varepsilon$. Under a stratified design with $m_g$
records from a stratum with machine share $\varepsilon_g$, $\E[K]=\sum_gm_g\varepsilon_g$,
which equals $m\varepsilon$ under proportional allocation.
\end{proposition}
For $N=900$, $m=45$ and $\varepsilon=0.05$, $\Pr(K\ge 4)=0.18$ and $\Pr(K\ge 6)=0.02$:
at low contamination a small audit rarely contains enough machine records, whatever
estimator is used afterwards. Only oversampling high-score strata raises $\E[K]$, and the
next result shows what that costs.

\begin{proposition}[Allocation]\label{p:neyman}
Under stratified simple random sampling with $\sum_gm_g=m$, ignoring integer constraints
and $m_g\le N_g$, the design variance of $\hat t_e$ is
$\sum_gN_g^2\sigma_g^2/m_g-\sum_gN_g\sigma_g^2$, where $\sigma_g^2$ is the variance of $e$ in
stratum $g$ (divisor $N_g-1$). It is minimised by $m_g\propto N_g\sigma_g$, with minimum
$(\sum_gN_g\sigma_g)^2/m-\sum_gN_g\sigma_g^2$; proportional allocation exceeds the minimum
by $(N/m)\sum_gN_g(\sigma_g-\bar\sigma)^2$ with $\bar\sigma=\sum_gN_g\sigma_g/N$.
\end{proposition}
\begin{proof}
By the Cauchy--Schwarz inequality
$(\sum_gN_g\sigma_g)^2\le(\sum_gN_g^2\sigma_g^2/m_g)(\sum_gm_g)$, with equality if and only if
$m_g\propto N_g\sigma_g$. Under proportional allocation the first term is
$(N/m)\sum_gN_g\sigma_g^2$, and
$\sum_gN_g\sigma_g^2-(\sum_gN_g\sigma_g)^2/N=\sum_gN_g(\sigma_g-\bar\sigma)^2$
(Cochran, 1977, Section~5.5).
\end{proof}
The optimal allocation follows the spread of the residual, not the score. Under the
working model of Remark~\ref{r:interval}, the residual spread of a stratum is largest where
the origin is uncertain ($p_g$ near one half) and outcomes are dispersed; a top score
decile in which nearly every record is machine-written and correctly predicted is not such
a place. Because $\sigma_g$ depends on unverified origins, an allocation close to Neyman's
must be learned from a pilot audit, and then the following applies.

\begin{lemma}[Two-stage audits]\label{l:twostage}
Let a pilot $S_1$ be drawn by stratified simple random sampling with sizes
$n_{1g}<N_g$, let the second-stage sizes $m_{2g}\ge1$ be any function of $S_1$ and of the
origins verified in it, and let $S_2$ be a simple random sample of $m_{2g}$ of the
$N_g-n_{1g}$ records of stratum $g$ outside $S_1$. For any $z$ not depending on the samples,
\begin{equation}\label{eq:twostage}
\hat t^{(2)}_z=\sum_{i\in S_1}z_i+\sum_g\frac{N_g-n_{1g}}{m_{2g}}\sum_{i\in S_2\cap g}z_i
\end{equation}
satisfies $\E[\hat t^{(2)}_z\mid S_1]=t_z$. Consequently:
(a) $\hat t^{(2)}_z$ is unbiased and $\Var(\hat t^{(2)}_z)=\E[\Var(\hat t^{(2)}_z\mid S_1)]$, so
the stratified variance estimator applied to the second stage alone is unbiased for it when
$m_{2g}\ge2$;
(b) the single-stage weights $N_g/(n_{1g}+m_{2g})$ applied to $S_1\cup S_2$ are in general
biased when $m_{2g}$ depends on the pilot origins;
(c) when the sizes $m_{2g}$ are fixed in advance, the variance of the stratum-$g$ term of
$\hat t^{(2)}_z$ is $(N_g-n_{1g})(n_{1g}+m_{2g})/(N_gm_{2g})$ times that of a single-stage
simple random sample of $n_{1g}+m_{2g}$ records.
\end{lemma}
\begin{proof}
Given $S_1$, the sizes $m_{2g}$ are fixed and the second sum is the stratified expansion
estimator of $\sum_{i\notin S_1}z_i$ on the population $U\setminus S_1$; (a) follows since
the conditional mean is the constant $t_z$. For (b) take one stratum with $N=3$,
$z=a=(1,0,0)$, $n_1=1$, and $m_2=1$ if the pilot record is machine-written and $m_2=2$
otherwise: the single-stage estimator has expectation
$\tfrac13\cdot\tfrac32\cdot1+\tfrac23\cdot1\cdot1=\tfrac76\neq1$. For (c), with $\sigma^2$ the
stratum variance of $z$ (divisor $N-1$) and $R$ the $N-n_1$ records outside the pilot, the conditional
variance is $(N-n_1)(N-n_1-m_2)\sigma_R^2/m_2$ and $\E[\sigma_R^2]=\sigma^2$ because $R$ is
itself a simple random sample; a single-stage sample of $n_1+m_2$ has variance
$N(N-n_1-m_2)\sigma^2/(n_1+m_2)$.
\end{proof}
Part (c) is the price of adaptivity: pilot records count only for themselves, and with a
pilot of 30\% of the budget the factor is about $1/0.7=1.4$ before any gain from a better
allocation. Whether a label-adaptive design nevertheless beats proportional allocation,
whose weights are known in advance and whose loss relative to Neyman allocation is given
in Proposition~\ref{p:neyman}, is an empirical question; at the budgets and with the
allocation rules studied here it did not on the development corpora, and the comparison
is repeated on the held-out corpora as a secondary analysis.

\subsection{Certification, verification error, and regret}\label{sec:cert}
Let $b=\bar y_{d=0}-\muH$ be the deletion bias on the response set and
$\hat b=\bar y_{d=0}-\hat\mu$ its estimate from an audit-based $\hat\mu$ with standard error
$\widehat{\mathrm{se}}$. Since $\bar y_{d=0}$ is computed from the whole response set,
$\hat b-b=-(\hat\mu-\muH)$.

\begin{proposition}[Certifying deletion]\label{p:tost}
Fix a tolerance $\tau>0$ and certify deletion when
$|\hat b|+q_{1-\alpha}\,\widehat{\mathrm{se}}<\tau$, with $q_{1-\alpha}$ at least the
standard normal $(1-\alpha)$ quantile. If $(\hat\mu-\muH)/\widehat{\mathrm{se}}$ is
asymptotically standard normal, then $\limsup\Pr(\text{certify})\le\alpha$ whenever
$|b|\ge\tau$.
\end{proposition}
\begin{proof}
If $b\ge\tau$, certification implies $\hat b+q_{1-\alpha}\widehat{\mathrm{se}}<\tau\le b$, that
is $(\hat\mu-\muH)/\widehat{\mathrm{se}}>q_{1-\alpha}$, an event of limiting probability at most
$\alpha$. The case $b\le-\tau$ is symmetric.
\end{proof}
This is the two one-sided tests procedure for equivalence. We use $\tau=0.2$ standard
deviations of the outcome in the response set, a substantive choice, and a Student
quantile with the degrees of freedom of Remark~\ref{r:interval}. The audit-based estimate
is the complete rule described after Theorem~\ref{t:ident} (RC-PPI++, or flag-PPI when
RC-PPI++ is infeasible); the normal limit is stated above, under the assumptions of
Theorem~\ref{t:clt}, for each estimator separately but not for the rule that switches between them on an origin-dependent event,
so the guarantee is asymptotic and conditional on that limit, and the false-certification
rate in small audits is measured (endpoint PE3). The rule of the first submission compared
a bootstrap interval of $\hat b$ with zero and had no such guarantee.

\begin{proposition}[Verification error]\label{p:flip}
Let origins be binary and suppose the audit returns $\tilde a_i$ instead of $a_i$, with
independent errors $\Pr(\tilde a_i=1\mid a_i=0)=\varphi_0<1$ and
$\Pr(\tilde a_i=0\mid a_i=1)=\varphi_1$ that do not depend on $y_i$ or on the design. For
fixed $(\eta,\lambda)$, $\E[\hat T_\lambda]/\E[\hat H_\lambda]=\tilde\mu$ with
\begin{equation}\label{eq:flip}
\tilde\mu-\muH=\frac{\varphi_1\,\varepsilon\,(\mu_A-\muH)}{(1-\varphi_0)(1-\varepsilon)+\varphi_1\varepsilon}.
\end{equation}
\end{proposition}
\begin{proof}
$\E[1-\tilde a_i]=(1-\varphi_0)h_i+\varphi_1a_i$. Repeating part (1) of
Theorem~\ref{t:ident} with this in place of $h_i$ gives
$\E[\hat T_\lambda]=(1-\varphi_0)N_H\muH+\varphi_1N_A\mu_A$ and
$\E[\hat H_\lambda]=(1-\varphi_0)N_H+\varphi_1N_A$.
\end{proof}
The quantity $\tilde\mu$ is the ratio of the expectations, to which the estimator is close
in large audits. Without machine records passed as human ($\varphi_1=0$) it equals $\muH$: humans wrongly
rejected ($\varphi_0$) alone cost precision, not validity, although they amplify the
shift when $\varphi_1>0$. Its distance from $\muH$ is at most
$\varphi_1\frac{\varepsilon}{1-\varepsilon}|\mu_A-\muH|/(1-\varphi_0)$, which gives a
sensitivity bound for stated bounds on $\varphi_0$, $\varphi_1$, $\varepsilon$ and the
range of the outcome; these cannot be estimated from the erroneous audit itself.

\begin{definition}[Regret]\label{d:regret}
A \emph{cell} is a corpus, detector, contamination rate, audit budget and outcome. For an
estimation rule $r$ let $\mathrm{sRMSE}_r(c)$ be its root mean squared error in cell $c$
over replications, divided by the standard deviation of the human outcome. The regret of
$r$ in $c$ is $\mathrm{sRMSE}_r(c)-\min_{r'\in\mathcal R}\mathrm{sRMSE}_{r'}(c)$, where
$\mathcal R$ is the set of single estimators (naive, deletion, two-phase H\'ajek, flag-PPI,
RC-PPI++). The comparator is the best single estimator in hindsight, which no analyst can
select without the truth; regret is non-negative for single estimators and may be negative
for a rule that switches between them.
\end{definition}

\begin{remark}[Selection policy]\label{r:policy}
The policy reports the candidate (naive, deletion, or the audit-based estimate just
described) with the smallest estimated mean squared error. For candidate
$j\in\{\text{naive},\text{deletion}\}$ the error is its bias $b_j$ on the response set,
estimated by $\hat b_j=\hat\mu_j-\hat\mu$; since $\E[\hat b_j^2]\approx b_j^2+\Var(\hat\mu)$,
the plug-in $\hat b_j^2$ overstates the squared bias and the policy uses
$\max(\hat b_j^2-\widehat{\mathrm{se}}^{\,2},0)$, compared with $\widehat{\mathrm{se}}^{\,2}$
for the audit-based estimate. Equivalently, naive or deletion is preferred to the
audit-based estimate only if $|\hat b_j|<\sqrt2\,\widehat{\mathrm{se}}$. If both have estimated squared bias zero, the naive mean is reported. The interval
reported with the policy is that of the audit-based estimate. Note the two scales: the
certification tolerance is in standard deviations of the outcome over all records of the
response set, which the analyst observes, whereas regret is standardised by the human
standard deviation, which only the evaluation knows.
\end{remark}
%%% END BODY

\section*{Numerical checks}
{\small\begin{tabular}{@{}lp{10.2cm}@{}}\toprule
Statement & Test (repository \texttt{tests/theory/})\\\midrule
Lemma~\ref{l:del}, Proposition~\ref{p:momentbias} & \texttt{test\_identities.py}: exact to $10^{-10}$ on response sets\\
Proposition~\ref{p:moment} & \texttt{test\_exact\_statements.py}: the solution equals $\muH$ on a response set built to satisfy M1 exactly\\
Theorem~\ref{t:ident} & \texttt{test\_exact\_statements.py}: unbiasedness by enumeration (simple and stratified sampling, graded origins, $\lambda$ outside $[0,1]$); representation \eqref{eq:repr} exact for the package estimators\\
Theorem~\ref{t:clt} & \texttt{test\_exact\_statements.py}: coverage of the linearised normal interval \eqref{eq:se} in a large audit (simulation); \texttt{test\_design\_based.py}: reported interval\\
Corollary~\ref{c:tune} & \texttt{test\_exact\_statements.py}: estimated variance at the package's $\hat\lambda$ never above that at $\lambda=0$, stratified audits\\
Remark~\ref{r:interval} & \texttt{test\_exact\_statements.py}: \eqref{eq:anticipated} by enumeration of all label vectors\\
Proposition~\ref{p:starve} & \texttt{test\_exact\_statements.py}: both probabilities\\
Proposition~\ref{p:neyman} & \texttt{test\_exact\_statements.py}: minimum against random allocations; excess of proportional allocation\\
Lemma~\ref{l:twostage} & \texttt{test\_exact\_statements.py}: (a)--(c) by enumeration; \texttt{tests/unit/test\_audit\_designs.py}: unbiased totals for the package's two-stage designs (simulation)\\
Proposition~\ref{p:tost} & \texttt{test\_exact\_statements.py}: certification rate at the boundary $|b|=\tau$ not above $\alpha$ (simulation, one-sided 1\% test)\\
Proposition~\ref{p:flip} & \texttt{test\_exact\_statements.py}: \eqref{eq:flip} exact; \texttt{test\_design\_based.py}: simulation\\
Remark~\ref{r:policy} & \texttt{test\_exact\_statements.py}: the package policy follows the $\sqrt2$ rule\\\bottomrule
\end{tabular}}
\end{document}
